%%%
%
% $Autor: Siddhanth Sharma $
% $Date: 2026-01-24 $
% $File: Monitoring.tex $
% $Version: 1.0 $
%
% !TeX spellcheck = en_GB
% !TeX program = pdflatex
% !TeX encoding = utf8
%
%%%

\chapter{Monitoring}
\label{ch:monitoring}
\index{Monitoring}

\noindent
This chapter documents how the running license plate recognition system is observed, validated, and maintained over time on an edge device (NVIDIA Jetson Nano).
The goal is a realistic, student-project scale approach: lightweight monitoring that fits the compute/storage constraints and works without cloud dependencies.

% --------------------------------------------------------
\section{Idea and Motivation}
\index{Monitoring!Idea and Motivation}

\noindent
Once deployed, the pipeline (camera $\rightarrow$ detection $\rightarrow$ OCR $\rightarrow$ database) can degrade due to environmental changes (lighting, camera angle), hardware limits (thermal throttling), or simple runtime issues (camera disconnects, corrupted frames).
Monitoring is therefore required to (i) detect failures early, (ii) preserve evidence for debugging, and (iii) collect new data for iterative improvement.

\noindent
The core objectives for this project are:
\begin{itemize}
	\item \textbf{Performance:} ensure the system stays responsive (e.g., stable processing rate / acceptable latency).
	\item \textbf{Hardware health:} detect resource pressure and thermal throttling on the Jetson Nano.
	\item \textbf{Availability:} detect crashes, camera failures, and model-load failures quickly.
	\item \textbf{Data quality:} identify unreadable/low-quality plate crops and collect them as new training data.
	\item \textbf{Operational correctness:} ensure that entry/exit logic and price calculation behave plausibly during a demo run (e.g., no duplicate charges for the same plate in a short time window).
\end{itemize}

% --------------------------------------------------------
\section{Plan / Description (Implementation Approach)}
\index{Monitoring!Plan / Description}

\subsection{What is monitored}
\index{Monitoring!Plan / Description!What is monitored}

\noindent
The project uses a small set of practical signals that are easy to collect locally:
\begin{itemize}
	\item \textbf{Detection/OCR quality:} plate text, OCR confidence, and whether the output matches a German plate format.
	\item \textbf{Throughput indicator:} processing interval / FPS-style feedback during testing (e.g., camera test scripts overlay FPS).
	\item \textbf{System resources (Jetson):} temperature, CPU/GPU load, memory usage, and disk space using built-in OS tools.
	\item \textbf{Operational events:} errors such as camera not opening, empty frames, and model-load failures.
	\item \textbf{Demo state signals:} displayed plate state (e.g., the overlay \texttt{READ: ...}), the live web UI table, and the computed fee from the SQLite state (implemented in \FILE{Code/LicencePlateDetection/jetson-inference/python/www/flask/dashboard2.py}).
\end{itemize}

\subsection{How metrics are collected and stored}
\index{Monitoring!Plan / Description!Collection and storage}

\noindent
To remain feasible on a Jetson Nano, metrics are stored locally and kept simple:
\begin{itemize}
	\item \textbf{Event log (SQLite):}
	the deployed dashboard persists a minimal \texttt{visits} table using \PYTHON{sqlite3} (\FILE{Code/LicencePlateDetection/jetson-inference/python/www/flask/dashboard2.py}).
	Each plate acts as a primary key and is tracked with an \textbf{entry timestamp} and a \textbf{last-seen timestamp}.
	This supports lightweight monitoring of \emph{currently seen cars} and \emph{demo billing state} via the local database file configured as \PYTHON{DB\_PATH} (default:\linebreak
	\texttt{/media/myssd/jetson-inference/python/www/flask/parking.db}).
	\item \textbf{Optional evidence capture (images):}
	the current deployed dashboard does \textbf{not} automatically persist plate crops for unreadable cases.
	For dataset growth, captures are obtained either by temporarily enabling \PYTHON{cv2.imwrite(...)} in \FILE{Code/LicencePlateDetection/jetson-inference/python/www/flask/dashboard2.py} (store the stable crop used for OCR), or by using a dedicated capture script during development (example: \FILE{Code/LicencePlateDetection/jetson\_app/detector.py} writes periodic samples to \FILE{Code/LicencePlateDetection/jetson\_app/captures}).
	\item \textbf{Resource snapshots (manual / scheduled):}
	during testing and demos, Jetson resource metrics can be observed with \SHELL{tegrastats} or by reading thermal sensors from \FILE{/sys/devices/virtual/thermal/}.
	If required, these outputs can be appended to a plain log file at a low frequency (e.g., every 10--30 seconds) to avoid overhead.
\end{itemize}

\subsection{Simple monitoring architecture}
\index{Monitoring!Plan / Description!Architecture}

\noindent
The monitoring concept follows the deployed data flow:
	\begin{enumerate}
		\item \textbf{Camera input} is validated (camera opens, frames are non-empty).
		\item \textbf{Detection} runs through a TensorRT engine and only processes boxes above a confidence threshold (configured via \PYTHON{CONFIDENCE\_THRESHOLD} in\linebreak
		\FILE{Code/LicencePlateDetection/jetson-inference/python/www/flask/dashboard2.py}).
		\item \textbf{OCR gating} is time-based: OCR is executed only after the bounding box is stable for at least \PYTHON{OCR\_DELAY} and not more frequently than \PYTHON{OCR\_COOLDOWN} (\FILE{Code/LicencePlateDetection/jetson-inference/python/www/flask/dashboard2.py}).
		\item \textbf{Normalization/validation} cleans OCR output into a German plate format using \PYTHON{clean\_german\_plate()} and a small corrections dictionary (\FILE{Code/LicencePlateDetection/jetson-inference/python/www/flask/dashboard2.py}).
		\item \textbf{Logging layer} writes a minimal state record (SQLite) by inserting/updating \texttt{visits} (plate $\rightarrow$ entry\_ts, last\_seen\_ts).
		\item \textbf{Operator checks} (during operation) confirm system resources and review logs/samples.
	\end{enumerate}

% --------------------------------------------------------
\section{Getting New Data During Operation}
\index{Monitoring!Getting New Data}

\noindent
To improve robustness over time, the system collects new samples from real operation instead of relying purely on the original dataset.
The collection strategy is intentionally selective to control storage usage on the edge device.

\subsection{Trigger-based capture}
\index{Monitoring!Getting New Data!Trigger-based capture}

\noindent
New data is captured under conditions that indicate that the current model/pipeline struggles:
\begin{itemize}
	\item \textbf{OCR fails / invalid format:} if the dashboard produces no valid plate string after
	\PYTHON{clean\_german\allowbreak\_plate()},
	the recommended workflow is to \emph{temporarily} enable crop saving in
	\FILE{Code/LicencePlateDetection/\allowbreak jetson-inference/\allowbreak python/www/flask/\allowbreak dashboard2.py}
	so that hard cases can be reviewed offline.
	
	\item \textbf{High-confidence detections:} for planned dataset collection runs, the development detector can
	periodically save frames and plate crops when detection confidence exceeds a threshold (example:
	\PYTHON{captureConfidence=0.5} with a minimum interval of \PYTHON{1500ms} in
	\FILE{Code/LicencePlateDetection/\allowbreak jetson\_app/\allowbreak detector.py}).
\end{itemize}
\noindent
This approach creates a focused "hard cases" dataset (blur, glare, angle issues, partial occlusions) that is valuable for retraining and debugging.

% --------------------------------------------------------
\section{Data Updating in the ML Pipeline}
\index{Monitoring!Data Updating in the ML Pipeline}

\noindent
Retraining in this student project is intentionally \textbf{offline and partially manual}.
The Jetson Nano focuses on inference, while dataset curation and training are performed on a development machine or a controlled environment.

\subsection{Review and annotation workflow}
\index{Monitoring!Data Updating in the ML Pipeline!Review and annotation}

\noindent
The stored captures are periodically reviewed:
\begin{itemize}
	\item Remove corrupted or unusable images (e.g., empty crops).
	\item Label new plates (bounding boxes in YOLO format) and, if required, create corresponding OCR ground truth for evaluation.
	\item Track the data source and conditions (night/day, rain, camera position) as lightweight metadata (directory naming or CSV).
\end{itemize}

\subsection{Retraining and model update}
\index{Monitoring!Data Updating in the ML Pipeline!Retraining and deployment}

\noindent
The model update strategy is:
\begin{enumerate}
	\item Train or fine-tune the detector offline (example script: \FILE{Code/LicencePlateDetection/scripts/train\_german\_lpr.py}).
	\item Validate the new model on a fixed local test set (same image set each iteration) and on recent "hard cases" collected from operation.
	\item Deploy by replacing the model artifact on the Jetson (e.g., updated \PYTHON{best.pt} weights) and re-running the startup checks.
\end{enumerate}

\noindent
Basic versioning is achieved by keeping dated model folders (e.g., \texttt{runs/train/...}) and preserving the previous working model for rollback.

% --------------------------------------------------------
\section{Checks (System Health Validation)}
\index{Monitoring!Checks}

\noindent
Checks are used at three times: \textbf{startup}, \textbf{runtime}, and \textbf{periodic review}.

\subsection{Startup checks}
\index{Monitoring!Checks!Startup}

\begin{itemize}
	\item \textbf{Python environment sanity:} verify imports and camera access using the project script
	\FILE{report/Code/SystemTools/\allowbreak python\_project\allowbreak\_sanity\allowbreak\_check.py}.
	
	\item \textbf{Engine availability:} verify the TensorRT engine file configured as
	\PYTHON{ENGINE\allowbreak\_PATH} exists before starting inference
	(\FILE{Code/LicencePlateDetection/\allowbreak jetson-inference/\allowbreak python/www/flask/\allowbreak dashboard2.py}).
	
	\item \textbf{Database readiness:} ensure the SQLite table \texttt{visits} exists on startup
	(\PYTHON{init\allowbreak\_db()} in
	\FILE{Code/LicencePlateDetection/\allowbreak jetson-inference/\allowbreak python/www/flask/\allowbreak dashboard2.py}).
	
	\item \textbf{Service supervision:} when started as a systemd service, verify \texttt{Restart=always}
	and inspect startup logs via \SHELL{systemctl status} / \SHELL{journalctl}
	(\FILE{Code/LicencePlateDetection/\allowbreak lpr-system.service}).
\end{itemize}

\subsection{Runtime checks}
\index{Monitoring!Checks!Runtime}

\begin{itemize}
	\item \textbf{Frame validity:} handle missing frames by continuing the loop and retrying
	(\FILE{Code/LicencePlateDetection/\allowbreak jetson-inference/\allowbreak python/www/flask/\allowbreak dashboard2.py}).
	
	\item \textbf{Camera fallback:} attempt camera index 0 and fall back to index 1 if needed
	(\FILE{Code/LicencePlateDetection/\allowbreak jetson-inference/\allowbreak python/www/flask/\allowbreak dashboard2.py}).
	
	\item \textbf{OCR throttling:} avoid spamming OCR using \PYTHON{OCR\allowbreak\_DELAY} and
	\PYTHON{OCR\allowbreak\_COOLDOWN} and require a minimum OCR confidence
	\PYTHON{OCR\allowbreak\_MIN\allowbreak\_CONFIDENCE}
	(\FILE{Code/LicencePlateDetection/\allowbreak jetson-inference/\allowbreak python/www/flask/\allowbreak dashboard2.py}).
	
	\item \textbf{Database writes:} handle and review DB write errors printed during insert/update transactions
	(\FILE{Code/LicencePlateDetection/\allowbreak jetson-inference/\allowbreak python/www/flask/\allowbreak dashboard2.py}).
\end{itemize}

\subsection{Periodic review checks}
\index{Monitoring!Checks!Periodic review}

\noindent
During longer runs or demos, logs and system resources are reviewed:
\begin{itemize}
	\item \textbf{SQLite growth:} ensure the database file does not grow without bounds:
	\texttt{/media/myssd/\allowbreak jetson-inference/\allowbreak python/www/flask/\allowbreak parking.db}
	(configured as \PYTHON{DB\allowbreak\_PATH} in
	\FILE{Code/LicencePlateDetection/\allowbreak jetson-inference/\allowbreak python/www/flask/\allowbreak dashboard2.py}).
	
	\item \textbf{Disk usage:} ensure the database partition and any optional capture directories do not fill the device storage.
	
	\item \textbf{Thermals:} check for thermal throttling symptoms (e.g., sudden FPS drops) and verify with Jetson OS tools.
	
	\item \textbf{Service logs:} review runtime errors and restarts in journald
	(service definition in \FILE{Code/LicencePlateDetection/\allowbreak lpr-system.service}).
\end{itemize}

% --------------------------------------------------------
\section{Code: Monitoring-Related Functions}
\index{Monitoring!Code Functions}

\noindent
To keep monitoring logic maintainable, the project separates core concerns into small functions/classes:
\begin{itemize}
	\item \textbf{Event persistence:} \PYTHON{init\allowbreak\_db()} creates and validates the SQLite table \texttt{visits},
	and the inference loop performs insert/update transactions
	(\FILE{Code/LicencePlateDetection/\allowbreak jetson-inference/\allowbreak python/www/flask/\allowbreak dashboard2.py}).
	
	\item \textbf{Plate validation:} \PYTHON{clean\allowbreak\_german\allowbreak\_plate()} normalizes and validates OCR output
	using a whitelist of city prefixes and simple formatting rules
	(\FILE{Code/LicencePlateDetection/\allowbreak jetson-inference/\allowbreak python/www/flask/\allowbreak dashboard2.py}).
	
	\item \textbf{Runtime UI/API monitoring:} \texttt{/api/billing} provides a direct view into the current database state,
	and the MJPEG stream endpoint exposes the latest processed frame
	(\FILE{Code/LicencePlateDetection/\allowbreak jetson-inference/\allowbreak python/www/flask/\allowbreak dashboard2.py}).
	
	\item \textbf{Performance feedback (testing):} \PYTHON{draw\allowbreak\_fps()} in
	\FILE{Code/LicencePlateDetection/\allowbreak scripts/\allowbreak test\allowbreak\_camera.py}
	provides quick runtime feedback during development.
\end{itemize}

% --------------------------------------------------------
\section{Code: Test Strategy}
\index{Monitoring!Code Test}

\noindent
Testing focuses on fast checks that catch deployment blockers early:
\begin{itemize}
	\item \textbf{Import and environment tests:} validate dependencies and basic runtime readiness (example: \FILE{Code/LicencePlateDetection/scripts/test\_yolo\_import.py} and \FILE{report/Code/SystemTools/python\_project\_sanity\_check.py}).
	\item \textbf{Camera and throughput tests:} validate camera read loop and visual overlays (example: \FILE{Code/LicencePlateDetection/scripts/test\_camera.py}).
	\item \textbf{End-to-end dry runs:} run the deployed application (\FILE{Code/LicencePlateDetection/jetson-inference/python/www/flask/dashboard2.py}) and confirm that (i) the web UI updates, (ii) \PYTHON{/api/billing} returns records, and (iii) the SQLite \texttt{visits} table is updated on valid plate reads.
\end{itemize}

% --------------------------------------------------------
\section{Privacy}
\index{Monitoring!Privacy}

\noindent
License plates are personal data in many jurisdictions, therefore monitoring must minimize data exposure:
\begin{itemize}
	\item \textbf{Data minimization:} the default operational log stores plate text and timestamps (SQLite) instead of storing full video streams or raw images.
	\item \textbf{Restricted access:} logs and captured images remain on the device and are intended for project team access only.
	\item \textbf{Retention:} captured images for debugging/retraining should be deleted after they have been reviewed and incorporated into an updated dataset; database entries should be pruned to a defined retention window (e.g., last N days) for the demo setup.
\end{itemize}

% --------------------------------------------------------
\section{Robustness}
\index{Monitoring!Robustness}

\noindent
The system is designed to degrade gracefully under common edge failures:
\begin{itemize}
	\item \textbf{Camera issues:} the application detects missing cameras and either falls back to another device index or terminates with a clear error.
	\item \textbf{OCR failures:} unreadable crops do not crash the system; they simply do not trigger a valid DB update (validity enforced via \PYTHON{clean\_german\_plate()} in \FILE{Code/LicencePlateDetection/jetson-inference/python/www/flask/dashboard2.py}).
	\item \textbf{GPU/CPU variability:} OCR is executed on the CPU (explicitly \PYTHON{gpu=False}) and can be quantified/optimized via EasyOCR settings (see initialization in\linebreak
	\FILE{Code/LicencePlateDetection/jetson-inference/python/www/flask/dashboard2.py}).
	\item \textbf{Automatic restart:} when installed as a systemd service, the application is restarted automatically on failure (\texttt{Restart=always} in \FILE{Code/LicencePlateDetection/lpr-system.service}).
\end{itemize}

% --------------------------------------------------------
\section{Process (Operational Workflow)}
\index{Monitoring!Process}

\noindent
A typical operational workflow for a demo run on the Jetson Nano is:
\begin{enumerate}
	\item \textbf{Start:} run the sanity check (\FILE{report/Code/SystemTools/python\_project\_sanity\_check.py}), then start the dashboard application.
	\item \textbf{Observe:} verify that the web UI updates, detections are shown in the stream, and that the SQLite database at \PYTHON{DB\_PATH} records visit state (see \FILE{Code/LicencePlateDetection/jetson-inference/python/www/flask/dashboard2.py}).
	\item \textbf{Monitor resources:} periodically check Jetson temperatures and utilization using OS tools during longer runs.
	\item \textbf{Stop:} terminate the application cleanly and archive logs/captures if needed.
	\item \textbf{Review:} if evidence capture is enabled (or if a dedicated capture run was performed), inspect failed samples and add useful images to the retraining dataset.
\end{enumerate}

\noindent
Example decision rules used during operation:
\begin{itemize}
	\item If the processing rate drops noticeably, check Jetson thermals and system load; reduce camera resolution or processing frequency if necessary.
	\item If many plates become unreadable, collect a small set of hard cases (temporary crop saving in \FILE{Code/LicencePlateDetection/jetson-inference/python/www/flask/dashboard2.py} or a dedicated capture run) to determine whether the cause is motion blur, lighting changes, or camera focus.
	\item If detections are missing entirely, run the startup checks again (camera access, model path, CUDA availability) before changing thresholds.
\end{itemize}
